\chapter{Page-by-Page Documentation}

Every substantive page in the application is documented below, in the
order a user would encounter them. Error pages (400/403/404/413/500)
are covered as a group at the end of this chapter rather than
individually, since they share one template pattern and no
business logic. Wireframes are simplified representations of the
actual rendered layout (sidebar + topbar are omitted from most
wireframes below the login page, since every authenticated page shares
the same shell --- see Chapter~6).

\section{Login}

\pf{Purpose}{The only entry point for an unauthenticated visitor;
establishes a session.}
\pf{Route}{\route{GET/POST /login}}
\pf{Access}{Public. An already-authenticated user is redirected to the
dashboard instead of seeing the form again.}
\pf{Main UI elements}{Email field, password field, submit button, a
flashed error banner on failure.}
\pf{Data displayed}{Nothing from the database --- a static form.}
\pf{User actions}{Submit credentials.}
\pf{Backend route}{\code{auth.login} (\file{app/routes/auth.py})}
\pf{Service functions}{\code{auth.service.authenticate}}
\pf{Database data involved}{\code{users} (read for the credential
check and lockout state; written on both success and failure to update
\code{failed\_login\_count}/\code{locked\_until}/\code{last\_login\_at}).}
\pf{Authorization}{None required to view; the route itself contains the
only authorization logic in the system that runs \emph{before} anyone is
authenticated.}
\pf{Validation}{Both fields required (WTForms \code{DataRequired}); no
email-format validation (deliberately --- an invalid format simply
won't match any account).}
\pf{Edge cases}{Wrong password, nonexistent email, inactive account,
and locked account all show the identical generic error (Chapter~2). A
successful login honors a same-site \route{?next=} redirect target and
rejects anything else.}
\pf{Error states}{A flashed banner; the form is always safe to
resubmit.}
\pf{Related pages}{Redirects successfully to the Dashboard.}
\pf{Source files}{\file{app/routes/auth.py}, \file{app/auth/service.py},
\file{templates/auth/login.html}}

\begin{center}
\begin{lstlisting}[basicstyle=\ttfamily\footnotesize,frame=single]
+----------------------------------------+
|      Workforce Management Platform      |
|                                          |
|   [ Email                            ]  |
|   [ Password                         ]  |
|              [ Log in ]                 |
+----------------------------------------+
\end{lstlisting}
\end{center}

\section{Dashboard --- Admin / Manager View}

\pf{Purpose}{"What is happening with the workforce right now" --- the
main operational landing page for admin/manager.}
\pf{Route}{\route{GET /dashboard}}
\pf{Access}{Admin (org-wide), manager (their managed departments).}
\pf{Main UI elements}{Metric grid (total employees, working today,
absent today, on leave today), a "Requiring Attention" panel, detail
tables (working/absent/on-leave/coverage), a labor-cost summary card,
an optional department filter.}
\pf{Data displayed}{Headline counts; a list of attendance entries
flagged \code{needs\_review} and employees missing pay-rate/policy
configuration; per-department coverage and labor-cost totals.}
\pf{User actions}{Filter by department; follow links into Attendance,
Employees, or Labor Cost for detail.}
\pf{Backend route}{\code{dashboard.index}, delegating to
\code{\_manager\_or\_admin\_dashboard} (\file{app/routes/dashboard.py})}
\pf{Service functions}{\code{reports.who\_is\_working\_today},
\code{reports.who\_is\_absent\_today},
\longcode{reports.who\_is\_on\_leave\_today},
\longcode{reports.attendance\_needing\_review},
\longcode{labor\_cost.department\_cost\_summary}}
\pf{Database data involved}{\code{shifts}, \code{attendance\_entries},
\code{leave\_requests}, \code{employees}, \code{employee\_pay\_rates},
\code{overtime\_policies} (indirectly, via labor cost).}
\pf{Authorization}{\code{role\_required} is not on this specific view
function --- \code{login\_required} covers the whole
\code{dashboard.index} route, which internally branches to the
employee dashboard for an employee-role caller; every underlying
\code{reports.*} call re-checks role/department scope itself.}
\pf{Validation}{An optional \route{?department\_id=} is a plain
integer filter; no format validation needed beyond that.}
\pf{Edge cases}{An employee with worked hours but no configured pay
rate/overtime policy no longer blanks the whole department's cost
total --- they are excluded and counted explicitly (\S4.6).}
\pf{Error states}{A \code{ValidationError} from the underlying labor
cost calculation is flashed, not raised as a 500.}
\pf{Related pages}{Employees, Attendance, Schedule, Leave, Labor Cost,
Reports --- this page is the hub linking to all of them.}
\pf{Source files}{\file{app/routes/dashboard.py},
\file{templates/dashboard/admin.html}}

\begin{center}
\begin{lstlisting}[basicstyle=\ttfamily\footnotesize,frame=single]
+------------------------------------------------------+
| Dashboard                                              |
+------------------------------------------------------+
| [Employees] [Working] [Absent] [On Leave]  <- metrics  |
+------------------------------------------------------+
| Requiring Attention (needs-review entries, gaps)       |
+------------------------------------------------------+
| Working Today   | Absent Today  | On Leave | Coverage  |
+------------------------------------------------------+
| Labor Cost (department totals)                         |
+------------------------------------------------------+
\end{lstlisting}
\end{center}

\section{Dashboard --- Employee View}

\pf{Purpose}{A personal "what's my status right now" landing page.}
\pf{Route}{\route{GET /dashboard} (same route, branches by role)}
\pf{Access}{Employee role only (own data).}
\pf{Main UI elements}{A prominent status card (Clock In / Currently
Working / Attendance Needs Review), a metric grid (upcoming shifts,
hours in the last 7 days, pending leave count), upcoming shifts table,
recent hours table, leave requests table.}
\pf{Data displayed}{The caller's own shifts, attendance, and leave
only.}
\pf{User actions}{Clock in / clock out; nothing else is writable from
this page.}
\pf{Backend route}{\code{dashboard.index} $\to$
\code{\_employee\_dashboard}}
\pf{Service functions}{\code{reports.current\_attendance\_status},
\code{scheduling.list\_shifts}, \longcode{working\_hours.scheduled\_vs\_worked},
\code{leave.list\_leave\_requests}}
\pf{Database data involved}{\code{shifts}, \code{attendance\_entries},
\code{leave\_requests} --- all filtered to \code{scope.employee\_id}.}
\pf{Authorization}{Every underlying call is scoped to the caller's own
\code{employee\_id}; there is no parameter on this page that could
target another employee.}
\pf{Validation}{None (no form inputs beyond the Clock In/Out buttons).}
\pf{Edge cases}{Three distinct states are rendered, not two: working,
\emph{needs review} (no button offered --- see \S4.4), or not clocked
in. This distinction exists specifically to avoid presenting a Clock
Out button that would always fail against a flagged entry.}
\pf{Error states}{A failed clock-in/out attempt (e.g.\ a race against
the open-entry index) is flashed as a friendly message, never a raw
500.}
\pf{Related pages}{Attendance (full history), Schedule (full upcoming
list), Leave.}
\pf{Source files}{\file{app/routes/dashboard.py},
\file{templates/dashboard/employee.html}}

\begin{center}
\begin{lstlisting}[basicstyle=\ttfamily\footnotesize,frame=single]
+------------------------------------------------------+
| My Dashboard                                           |
+------------------------------------------------------+
| [ Currently Working / Not Clocked In ]  [Clock Out/In] |
+------------------------------------------------------+
| [Upcoming Shifts] [Hours 7d] [Pending Leave]            |
+------------------------------------------------------+
| Upcoming Shifts table   | Recent Hours table            |
+------------------------------------------------------+
| My Leave Requests                                       |
+------------------------------------------------------+
\end{lstlisting}
\end{center}

\section{Employees --- List}

\pf{Purpose}{The organization's staff directory.}
\pf{Route}{\route{GET/POST /employees}}
\pf{Access}{Admin (all employees), manager (their departments). Not
reachable by an employee-role user at all.}
\pf{Main UI elements}{Department/status filters, an "Add employee"
disclosure (admin only), a table of employees with avatar-initial,
name, employee number, department, status.}
\pf{Data displayed}{Every employee visible to the caller's scope.}
\pf{User actions}{Filter; create a new employee (admin only); click
through to a detail page.}
\pf{Backend route}{\code{employees.list\_employees} /
\code{employees.create\_employee}}
\pf{Service functions}{\code{employees.list\_employees},
\code{employees.create\_employee}, \code{departments.list\_departments}}
\pf{Database data involved}{\code{employees}, \code{departments}.}
\pf{Authorization}{\code{role\_required("admin", "manager")} on the
route; \code{create\_employee} additionally checks
\code{scope.role == "admin"} inside the service.}
\pf{Validation}{Department/status filters are plain query parameters,
filtered in Python over an already-scoped list (no route-level
whitelist on the status value --- \VERIFY{} whether this should reject
an unrecognized status rather than silently returning no results).}
\pf{Edge cases}{A filtered view with zero matches is distinguished from
a genuinely empty roster ("No employees match this filter" vs.\ "No
employees yet").}
\pf{Error states}{An \code{IntegrityError} on create (duplicate
employee number/email) is flashed, not a raw 500.}
\pf{Related pages}{Employee Detail, Departments, Schedule (assignment).}
\pf{Source files}{\file{app/routes/employees.py},
\file{templates/employees/list.html}}

\begin{center}
\begin{lstlisting}[basicstyle=\ttfamily\footnotesize,frame=single]
+------------------------------------------------------+
| Employees                          [+ Add employee]   |
+------------------------------------------------------+
| Department [v]  Status [v]              [ Filter ]    |
+------------------------------------------------------+
| Employee          | Department | Status                |
| JS John Smith     | Kitchen    | Active                |
+------------------------------------------------------+
\end{lstlisting}
\end{center}

\section{Employee Detail}

\pf{Purpose}{A single employee's full profile, plus (for admin/manager)
a quick view of their schedule and attendance.}
\pf{Route}{\route{GET/POST /employees/<id>}}
\pf{Access}{Admin (any employee), manager (their departments), or the
employee themselves (their own record only).}
\pf{Main UI elements}{Profile card (contact info, status, linked login
account), an edit form (admin, or manager within scope), a terminate
form (admin only), and --- admin/manager only --- Upcoming Shifts and
Recent Attendance cards.}
\pf{Data displayed}{Employee profile fields; the linked \code{User}
account's email/role/active status, if any; (admin/manager) the next
14 days of shifts and the last 7 days of attendance.}
\pf{User actions}{Edit profile; terminate (admin); navigate to Pay Rate
or Labor Cost (admin only, via links in the page header).}
\pf{Backend route}{\code{employees.get\_employee} /
\code{employees.update\_employee} / \code{employees.terminate\_employee}}
\pf{Service functions}{\code{employees.get\_employee},
\code{employees.get\_linked\_user}, \code{scheduling.list\_shifts},
\code{reports.attendance\_entries\_with\_context}}
\pf{Database data involved}{\code{employees}, \code{users},
\code{shifts}, \code{attendance\_entries}.}
\pf{Authorization}{\code{employees.get\_employee} 404s an employee
attempting to view anyone but themselves; the Upcoming Shifts/Recent
Attendance sections are computed and rendered only when \code{can\_edit}
is true (admin, or manager within scope) --- an employee viewing their
own page never sees them (that information already lives on their own
dashboard).}
\pf{Validation}{Standard WTForms field validation on the edit form.}
\pf{Edge cases}{An employee with no linked login account shows "No
login account linked" rather than a blank field.}
\pf{Error states}{A 404 for an out-of-scope employee id (never a 403,
per the IDOR-defense convention in Chapter~2).}
\pf{Related pages}{Employees list, Pay Rate, Labor Cost breakdown,
Schedule, Attendance.}
\pf{Source files}{\file{app/routes/employees.py},
\file{templates/employees/detail.html}}

\begin{center}
\begin{lstlisting}[basicstyle=\ttfamily\footnotesize,frame=single]
+------------------------------------------------------+
| John Smith                    [Edit] [Pay rate] [Cost] |
+------------------------------------------------------+
| Profile: number, department, status, email, account    |
+------------------------------------------------------+
| Upcoming Shifts table   | Recent Attendance table       |
+------------------------------------------------------+
\end{lstlisting}
\end{center}

\section{Employee Pay Rate}

\pf{Purpose}{View and set an employee's hourly-rate history.}
\pf{Route}{\route{GET/POST /employees/<id>/pay-rate}}
\pf{Access}{Admin only --- never a manager, per rule A4.}
\pf{Main UI elements}{A history table (rate, effective from/to) and a
"set new rate" form.}
\pf{Data displayed}{Every historical rate period for the employee.}
\pf{User actions}{Add a new effective-dated rate.}
\pf{Backend route}{\code{employees.view\_pay\_rate} /
\code{employees.set\_pay\_rate}}
\pf{Service functions}{\code{pay\_rates.list\_pay\_rate\_history},
\code{pay\_rates.set\_pay\_rate}}
\pf{Database data involved}{\code{employee\_pay\_rates}.}
\pf{Authorization}{\code{role\_required("admin")} on the route,
re-checked inside both service functions.}
\pf{Validation}{Rate must be $> 0$; effective-to (if given) must be on
or after effective-from; the database's exclusion constraint is the
final authority against an overlapping period.}
\pf{Edge cases}{An overlapping rate period is rejected with a friendly
message, not a raw constraint-violation error.}
\pf{Error states}{\code{ValidationError} flashed on an invalid range
or overlap.}
\pf{Related pages}{Employee Detail, Labor Cost (employee breakdown).}
\pf{Source files}{\file{app/routes/employees.py},
\file{app/services/pay\_rates.py}, \file{templates/employees/pay\_rate.html}}

\section{Departments}

\pf{Purpose}{Manage the organization's department list.}
\pf{Route}{\route{GET/POST /departments}, \route{POST
/departments/<id>}, \route{POST /departments/<id>/deactivate}}
\pf{Access}{Admin (full CRUD + deactivate); manager (view only, scoped
to managed departments).}
\pf{Main UI elements}{A table of departments (name, code, status), an
edit form per row (admin), a create form (admin), a deactivate action
(admin).}
\pf{Data displayed}{Every department visible to the caller.}
\pf{User actions}{Create, edit, deactivate (admin only).}
\pf{Backend route}{\code{departments.list\_departments} /
\code{create\_department} / \code{update\_department} /
\code{deactivate\_department}}
\pf{Service functions}{\code{departments.*} (all of them).}
\pf{Database data involved}{\code{departments}.}
\pf{Authorization}{List route is \code{role\_required("admin",
"manager")}; every write route is \code{role\_required("admin")},
re-checked inside the service.}
\pf{Validation}{Name/code required; case-insensitive name uniqueness
and \code{(organization\_id, code)} uniqueness enforced at the
database level.}
\pf{Edge cases}{Deactivation never deletes the row --- employees and
historical schedules may still reference it.}
\pf{Error states}{A duplicate name/code on create/update is flashed as
"name or code already in use," not a raw constraint error.}
\pf{Related pages}{Employees, Schedule, Labor Cost.}
\pf{Source files}{\file{app/routes/departments.py},
\file{templates/departments/list.html}}

\section{Schedule}

\pf{Purpose}{The planned-work calendar.}
\pf{Route}{\route{GET/POST /schedule}, plus \route{/schedule/<id>},
\route{/schedule/<id>/assign}, \route{/schedule/<id>/publish},
\route{/schedule/<id>/cancel}}
\pf{Access}{Everyone can view (employees see only their own published
shifts); admin/manager can create/edit/assign/publish/cancel.}
\pf{Main UI elements}{Date-range and department filters, a "New shift"
disclosure, a table (date, start, end, hours, department, employee,
status, actions).}
\pf{Data displayed}{Shifts in the selected window, scoped by role.}
\pf{User actions}{Create, edit (draft only), assign/reassign an
employee, publish, cancel.}
\pf{Backend route}{\code{schedule.list\_shifts} and the write actions
listed above.}
\pf{Service functions}{\code{scheduling.*} (all of them).}
\pf{Database data involved}{\code{shifts}, \code{employees},
\code{departments}, and (for conflict checks) \code{leave\_requests}.}
\pf{Authorization}{View is \code{login\_required}; every write is
\code{role\_required("admin", "manager")}, with manager department
scoping re-checked inside every \code{scheduling.*} write function.}
\pf{Validation}{End time after start time and a 24-hour span cap are
database-enforced; overlap and approved-leave conflicts are checked
before every create/reassign/time-edit.}
\pf{Edge cases}{An unassigned shift shows a warning badge, not silent
blank text; a shift cannot be published without an assigned employee;
an out-of-range date filter shows an explicit "invalid date range"
state rather than a misleading "no shifts" message.}
\pf{Error states}{Overlap/leave-conflict violations are flashed as
friendly, specific messages (Chapter~3's exclusion constraints,
translated by \code{\_commit\_or\_raise\_overlap}).}
\pf{Related pages}{Employees, Leave (conflict source), Dashboard
(coverage summary).}
\pf{Source files}{\file{app/routes/schedule.py},
\file{templates/schedule/list.html}}

\begin{center}
\begin{lstlisting}[basicstyle=\ttfamily\footnotesize,frame=single]
+------------------------------------------------------+
| Schedule                              [+ New shift]   |
+------------------------------------------------------+
| From [ ] To [ ]  Department [v]         [ Filter ]     |
+------------------------------------------------------+
| Date  Start  End  Hours  Dept  Employee  Status  Acts  |
+------------------------------------------------------+
\end{lstlisting}
\end{center}

\section{Attendance}

\pf{Purpose}{The actual record of who worked when; the primary
clock-in/clock-out surface.}
\pf{Route}{\route{GET /attendance}, \route{POST /attendance/clock-in},
\route{POST /attendance/<id>/clock-out}, \route{POST
/attendance/<id>/correct}}
\pf{Access}{Everyone can view (employees see only their own entries);
admin/manager can correct any entry in scope and clock in/out on behalf
of others.}
\pf{Main UI elements}{A personal status card (employee only), a
date-range filter plus an employee filter (admin/manager), a table
(date, started, ended, status, worked hours, source, actions).}
\pf{Data displayed}{Attendance entries in the selected window, each
enriched with computed worked-duration and lateness against a matched
shift.}
\pf{User actions}{Clock in/out (self, or on behalf of others for
admin/manager); correct an entry (admin/manager, mandatory reason).}
\pf{Backend route}{\code{attendance.list\_entries} and the write
actions listed above.}
\pf{Service functions}{\longcode{attendance.clock\_in/clock\_out/correct\_entry},
\longcode{reports.attendance\_entries\_with\_context},
\longcode{reports.current\_attendance\_status}}
\pf{Database data involved}{\code{attendance\_entries}, \code{shifts}
(for the matched-shift/lateness display).}
\pf{Authorization}{View is \code{login\_required}; \code{correct\_entry}
is \code{role\_required("admin", "manager")}; clock-in/out on someone
else's behalf requires admin/manager, re-checked inside the service.}
\pf{Validation}{A backdated clock-in/out (admin/manager only) is capped
at 90 days in the past and can never be in the future; a correction
requires a non-blank reason.}
\pf{Edge cases}{A \code{needs\_review} entry never offers a Clock Out
button (it would always fail) --- only a correction can resolve it;
this status is detected by an unbounded lookup on
\code{ended\_at IS NULL}, not a recent-days window, specifically so an
old, still-unresolved entry is never missed (see the note on this fix
in Chapter~7).}
\pf{Error states}{A friendly message for a duplicate clock-in, an
already-closed entry, or a needs-review clock-out attempt --- never a
raw 500 from the underlying database constraint.}
\pf{Related pages}{Dashboard (personal status card), Employee Detail
(recent attendance), Schedule (matched shift source).}
\pf{Source files}{\file{app/routes/attendance.py},
\file{templates/attendance/list.html}}

\begin{center}
\begin{lstlisting}[basicstyle=\ttfamily\footnotesize,frame=single]
+------------------------------------------------------+
| Attendance                                             |
+------------------------------------------------------+
| [ Currently Working / Not Clocked In ]  [Clock Out/In] |
+------------------------------------------------------+
| From [ ] To [ ]  Employee [v]           [ Filter ]     |
+------------------------------------------------------+
| Date  Started  Ended  Status  Worked  Source  Actions  |
+------------------------------------------------------+
\end{lstlisting}
\end{center}

\section{Leave}

\pf{Purpose}{Request and manage time off.}
\pf{Route}{\route{GET/POST /leave}, plus approve/reject/cancel actions}
\pf{Access}{Everyone can view (employees see only their own requests)
and request leave; admin/manager can approve/reject/cancel within
scope.}
\pf{Main UI elements}{Status and employee filters (admin/manager), a
request form, a table (dates, employee, type, status, reason, actions),
inline conflict warnings.}
\pf{Data displayed}{Leave requests in scope; for a pending request,
any conflicting published shift.}
\pf{User actions}{Request leave; approve/reject (with a mandatory
reason for rejection); cancel.}
\pf{Backend route}{\code{leave.list\_requests} and the write actions.}
\pf{Service functions}{\code{leave.*} (all of them),
\code{leave.conflicting\_shifts\_for}}
\pf{Database data involved}{\code{leave\_requests}, \code{leave\_types},
\code{shifts} (conflict detection).}
\pf{Authorization}{View/request is \code{login\_required};
approve/reject is \code{role\_required("admin", "manager")}; cancel is
open to the owner (while pending) or admin/manager (any status),
enforced inside \code{cancel\_leave} itself.}
\pf{Validation}{Rejection requires a non-blank reason; an overlapping
pending/approved request for the same employee is rejected by the
database's exclusion constraint.}
\pf{Edge cases}{Self-approval is explicitly blocked, even for an
admin/manager who is also the requesting employee; approving over a
published shift is blocked with the specific conflicting shift(s)
named.}
\pf{Error states}{A friendly, specific conflict message rather than a
generic failure.}
\pf{Related pages}{Schedule (conflict source), Dashboard ("on leave
today").}
\pf{Source files}{\file{app/routes/leave.py},
\file{templates/leave/list.html}}

\section{Overtime Report}

\pf{Purpose}{Per-employee overtime hours for a department/date range.}
\pf{Route}{\route{GET /reports/overtime}}
\pf{Access}{Admin, manager (their departments).}
\pf{Main UI elements}{Department + date-range filter, a table
(employee, overtime hours), a total row, cross-navigation to Hours
Trend and Labor Cost for the same filter.}
\pf{Data displayed}{Overtime \emph{hours} only --- never a rate or a
cost figure (rule A4).}
\pf{User actions}{Filter; navigate to related reports.}
\pf{Backend route}{\code{dashboard.overtime\_report}}
\pf{Service functions}{\code{reports.overtime\_summary}}
\pf{Database data involved}{\code{attendance\_entries},
\code{overtime\_policies}/\code{overtime\_tiers}, \code{employee\_pay\_rates}
(read internally to price line items, never rendered here).}
\pf{Authorization}{\code{role\_required("admin", "manager")}, with
manager department scoping re-checked inside the service.}
\pf{Validation}{An excessively wide date range is clamped to 92 days
before the underlying per-employee/per-day calculation runs (Chapter~6).}
\pf{Edge cases}{An employee with no configured pay rate/policy shows
"Not available" for that row instead of failing the whole report.}
\pf{Error states}{None beyond the per-row "not available" case ---
the report always renders.}
\pf{Related pages}{Hours Trend, Labor Cost, Dashboard.}
\pf{Source files}{\file{app/routes/dashboard.py},
\file{templates/reports/overtime.html}}

\section{Hours Trend Report}

\pf{Purpose}{"Are working hours going up or down over time" for a
department.}
\pf{Route}{\route{GET /reports/hours-trend}}
\pf{Access}{Admin, manager (their departments).}
\pf{Main UI elements}{Department + date-range filter, a per-day bar
comparison of total worked hours, cross-navigation to Overtime and
Labor Cost.}
\pf{Data displayed}{Total worked hours per day, summed across the
department's employees.}
\pf{User actions}{Filter; navigate to related reports.}
\pf{Backend route}{\code{dashboard.hours\_trend\_report}}
\pf{Service functions}{\code{reports.hours\_trend}}
\pf{Database data involved}{\code{attendance\_entries}.}
\pf{Authorization}{\code{role\_required("admin", "manager")}.}
\pf{Validation}{Same 92-day clamp as the Overtime report; an
\code{end < start} range shows an explicit "invalid date range" state.}
\pf{Edge cases}{Bar width is computed server-side and mapped onto a
small set of CSS width classes (never an inline \code{style}
attribute, to respect the application's strict CSP --- see Chapter~6).}
\pf{Error states}{None beyond the invalid-range state.}
\pf{Related pages}{Overtime Report, Labor Cost, Dashboard.}
\pf{Source files}{\file{app/routes/dashboard.py},
\file{templates/reports/hours\_trend.html}}

\section{Labor Cost --- Department Totals}

\pf{Purpose}{Aggregate labor cost for a department/date range.}
\pf{Route}{\route{GET /labor-cost}}
\pf{Access}{Admin, manager (their departments).}
\pf{Main UI elements}{Department + date-range filter, a total-cost
card, an unconfigured-employee warning, cross-navigation to Hours
Trend/Overtime, and (admin only) a per-employee drill-down list linking
to individual breakdowns.}
\pf{Data displayed}{A single aggregate Decimal total --- never a rate,
never a per-employee figure, for a manager.}
\pf{User actions}{Filter; (admin) drill into one employee's breakdown.}
\pf{Backend route}{\code{labor\_cost.department\_totals}}
\pf{Service functions}{\code{labor\_cost.department\_cost\_summary}}
\pf{Database data involved}{\code{attendance\_entries},
\code{employee\_pay\_rates}, \code{overtime\_policies}/\code{tiers}
(internally).}
\pf{Authorization}{\code{role\_required("admin", "manager")}; the
per-employee drill-down section is computed and rendered only when
\code{is\_admin} --- a manager never sees even the \emph{existence} of
that link, not just its numbers.}
\pf{Validation}{Same 92-day range clamp as the other reports;
\code{end < start} raises a \code{ValidationError}, flashed rather than
crashing.}
\pf{Edge cases}{The per-employee drill-down list includes
\emph{every} employment status (not only active), matching exactly the
employee set the total above it was computed from --- otherwise the
total and the list of links would silently disagree about who
contributed to it.}
\pf{Error states}{A flashed message on an invalid range; a
"no cost total recorded" state when a valid department/range yields no
data.}
\pf{Related pages}{Overtime Report, Hours Trend, Labor Cost Employee
Detail (admin only).}
\pf{Source files}{\file{app/routes/labor\_cost.py},
\file{templates/labor\_cost/index.html}}

\section{Labor Cost --- Employee Breakdown}

\pf{Purpose}{The one view in the system that shows an individual
employee's hourly rate and exact cost, broken down by day and category.}
\pf{Route}{\route{GET /labor-cost/employees/<id>}}
\pf{Access}{Admin only --- structurally unreachable by a manager
(hard-coded \code{role\_required("admin")}, not
\code{role\_required("admin", "manager")} like almost every other
manager-adjacent view).}
\pf{Main UI elements}{Date-range filter, a line-item table (date,
category, hours, rate, multiplier, cost), a total row.}
\pf{Data displayed}{Every priced line item for the employee/range,
including the hourly rate in force on each day.}
\pf{User actions}{Filter by date range.}
\pf{Backend route}{\code{labor\_cost.employee\_detail}}
\pf{Service functions}{\code{labor\_cost.range\_cost\_for\_employee}}
\pf{Database data involved}{\code{attendance\_entries},
\code{employee\_pay\_rates}, \code{overtime\_policies}/\code{tiers}.}
\pf{Authorization}{\code{role\_required("admin")}, plus an org-scoped
404 lookup on the employee id (defense in depth, since the route is
already admin-only).}
\pf{Validation}{Same range clamp; a day with genuinely worked hours but
no configured rate/policy raises a specific
\code{ValidationError}, flashed rather than silently priced as zero.}
\pf{Edge cases}{A mid-week rate change can legitimately produce two
differently-priced weekly-overtime line items for the same tier
(\S4.6).}
\pf{Error states}{A flashed configuration-gap message; the table
simply shows whatever line items \emph{did} resolve.}
\pf{Related pages}{Employee Detail, Pay Rate, Labor Cost Department
Totals.}
\pf{Source files}{\file{app/routes/labor\_cost.py},
\file{templates/labor\_cost/employee\_detail.html}}

\section{Audit Log}

\pf{Purpose}{A complete, permanent, admin-only record of privileged
actions.}
\pf{Route}{\route{GET /audit-log}}
\pf{Access}{Admin only.}
\pf{Main UI elements}{A date-range filter, a paginated table (when,
action, entity, actor user id, IP address, details), Previous/Next
links.}
\pf{Data displayed}{Every audit entry in the selected window, 50 per
page.}
\pf{User actions}{Filter by date; page through results.}
\pf{Backend route}{\code{audit.list\_entries}}
\pf{Service functions}{\code{audit.list\_entries}}
\pf{Database data involved}{\code{audit\_logs}.}
\pf{Authorization}{\code{role\_required("admin")}, re-checked inside
\code{audit.list\_entries} itself (this route historically queried the
database directly; it now goes through the service like every other
route --- see Chapter~1's "Round C fix" note).}
\pf{Validation}{Page number is clamped to a minimum of 1; an invalid
date string falls back to the default window rather than erroring.}
\pf{Edge cases}{The actor's raw user id is shown, not a resolved email
--- unlike the Recent Activity view below --- \VERIFY{} whether this
asymmetry is intentional or an omission.}
\pf{Error states}{An empty range shows a clear "no audit events in
this range" state.}
\pf{Related pages}{Recent Activity (a lighter view over the same
data).}
\pf{Source files}{\file{app/routes/audit.py},
\file{app/services/audit.py}, \file{templates/audit/list.html}}

\section{Recent Activity}

\pf{Purpose}{A lighter, easier-to-read glance at the most recent
privileged actions, for a quick check rather than a full
investigation.}
\pf{Route}{\route{GET /recent-activity}}
\pf{Access}{Admin only.}
\pf{Main UI elements}{A date-range filter, an activity list with each
entry's actor email resolved, a link through to the full Audit Log.}
\pf{Data displayed}{The same underlying audit events as the Audit Log,
capped at \code{audit.DEFAULT\_PAGE\_SIZE} (50) with no pagination
controls.}
\pf{User actions}{Filter by date; follow through to the full log.}
\pf{Backend route}{\code{audit.recent\_activity}}
\pf{Service functions}{\code{reports.recent\_activity} (which itself
calls \code{audit.list\_entries} and resolves actor emails in one
batched query).}
\pf{Database data involved}{\code{audit\_logs}, \code{users} (actor
email resolution).}
\pf{Authorization}{\code{role\_required("admin")}, re-checked inside
\code{reports.recent\_activity} itself.}
\pf{Validation}{Same date-parsing fallback as the Audit Log.}
\pf{Edge cases}{An entry whose actor account no longer resolves (should
not normally happen, since users are never hard-deleted) shows no
email rather than erroring.}
\pf{Error states}{Same "no events" empty state pattern as the Audit
Log.}
\pf{Related pages}{Audit Log.}
\pf{Source files}{\file{app/routes/audit.py},
\file{app/services/reports.py}, \file{templates/audit/recent\_activity.html}}

\section{Error Pages}

Five static error templates exist, one per handled HTTP status:
\file{errors/400.html}, \file{403.html}, \file{404.html},
\file{413.html} (payload too large --- reachable because
\code{MAX\_CONTENT\_LENGTH} is set to 1\,MB), and \file{500.html}. None
of them ever renders a stack trace, a SQL error, or any other internal
detail --- \code{app.errors.handle\_internal\_server\_error} logs the
real exception server-side and returns only a generic message. None
require authorization of their own (a 404 is deliberately returned
\emph{instead of} a 403 for an out-of-scope resource, per Chapter~2's
IDOR defense, so these pages are reached by every role equally).

\section{Admin Dashboard Deep Dive}

The Admin Dashboard (\S5.2) is the most information-dense page in the
application, so its data flow is documented separately here.

\begin{center}
\resizebox{0.95\textwidth}{!}{%
\begin{tikzpicture}[node distance=0.6cm]
  \node[flownode, text width=3.4cm] (shifts) {Shift data};
  \node[flownode, text width=3.4cm, below=of shifts] (att) {Attendance data};
  \node[flownode, text width=3.4cm, below=of att] (leave) {Leave data};
  \node[flownode, text width=3.4cm, below=of leave] (cost) {Labor cost data};

  \node[flownode, text width=5.2cm, right=2.2cm of att] (svc) {\code{reports.*} / \code{labor\_cost.*}\\(dashboard service composition)};
  \node[flownode, text width=5.2cm, below=1.0cm of svc] (route) {\code{dashboard.\_manager\_or\_admin\_dashboard}};
  \node[flownode, text width=5.2cm, below=1.0cm of route] (tmpl) {\file{dashboard/admin.html} (Jinja)};

  \draw[flowarrow] (shifts) -- (svc);
  \draw[flowarrow] (att) -- (svc);
  \draw[flowarrow] (leave) -- (svc);
  \draw[flowarrow] (cost) -- (svc);
  \draw[flowarrow] (svc) -- (route);
  \draw[flowarrow] (route) -- (tmpl);
\end{tikzpicture}%
}
\end{center}

\textbf{Metrics.} Total employees, currently working, absent, on leave
--- each a thin composition over an existing \code{reports.who\_is\_*}
function, not a new query.

\textbf{Workforce status / attendance.} \code{who\_is\_working\_today}
and \code{who\_is\_absent\_today} together define "absent" as
\emph{scheduled but with no attendance entry at all today, and not on
approved leave} --- not "did not clock into the specific shift they
were assigned," a distinction that matters for an employee with two
shifts or an early clock-in outside the shift-matching grace window
(documented in \code{reports.py}'s "Round B fix" comment).

\textbf{Departments / coverage / labor cost.} Computed once per visible
department via \code{scheduling.coverage\_summary} and
\code{labor\_cost.department\_cost\_summary}.

\textbf{Recent activity.} Available as its own page (\S5.15), reached
from the sidebar, not embedded on the dashboard itself.

\textbf{Filters.} An optional department filter narrows every section
above simultaneously.

\textbf{Admin-only vs.\ manager-safe information.} Everything on this
page is manager-safe \emph{except} that a manager only ever sees their
own managed departments' data (the query scoping, not a separate
rendering path, is what enforces this --- there is no
"admin-only widget" hidden from a manager on this particular page,
unlike the Labor Cost and Pay Rate pages).
